\documentclass[12pt,a4paper]{article}
\usepackage[utf8]{inputenc}
\usepackage[english,french]{babel}
\usepackage[T1]{fontenc}
\usepackage{amsmath,amsfonts,amssymb}
\usepackage{graphicx}
\usepackage{xcolor}
\usepackage{listings}
\usepackage{listingsutf8}
\usepackage{geometry}
\usepackage{hyperref}
\usepackage{titlesec}
\usepackage{fancyhdr}
\usepackage{tcolorbox}
\usepackage{enumitem}

\geometry{margin=2.5cm}
\setlength{\headheight}{15pt}
\hypersetup{
    colorlinks=true,
    linkcolor=blue,
    filecolor=magenta,      
    urlcolor=cyan,
    pdftitle={LangChain4j Integration Guide},
    pdfauthor={AI System Architecture},
    pdfsubject={Multi-Agent RAG System with Spring Boot}
}

% Custom colors
\definecolor{codegreen}{rgb}{0,0.6,0}
\definecolor{codegray}{rgb}{0.5,0.5,0.5}
\definecolor{codepurple}{rgb}{0.58,0,0.82}
\definecolor{backcolour}{rgb}{0.95,0.95,0.92}

% Code styling
\lstdefinestyle{mystyle}{
    backgroundcolor=\color{backcolour},   
    commentstyle=\color{codegreen},
    keywordstyle=\color{magenta},
    numberstyle=\tiny\color{codegray},
    stringstyle=\color{codepurple},
    basicstyle=\ttfamily\footnotesize,
    breakatwhitespace=false,         
    breaklines=true,                 
    captionpos=b,                    
    keepspaces=true,                 
    numbers=left,                    
    numbersep=5pt,                  
    showspaces=false,                
    showstringspaces=false,
    showtabs=false,                  
    tabsize=2
}
\lstset{style=mystyle}

% Define JSON language for listings
\lstdefinelanguage{json}{
    basicstyle=\normalfont\ttfamily,
    numbers=left,
    numberstyle=\scriptsize,
    stepnumber=1,
    numbersep=8pt,
    showstringspaces=false,
    breaklines=true,
    frame=lines,
    backgroundcolor=\color{backcolour},
    literate=
     *{0}{{{\color{numb}0}}}{1}
      {1}{{{\color{numb}1}}}{1}
      {2}{{{\color{numb}2}}}{1}
      {3}{{{\color{numb}3}}}{1}
      {4}{{{\color{numb}4}}}{1}
      {5}{{{\color{numb}5}}}{1}
      {6}{{{\color{numb}6}}}{1}
      {7}{{{\color{numb}7}}}{1}
      {8}{{{\color{numb}8}}}{1}
      {9}{{{\color{numb}9}}}{1}
      {:}{{{\color{punct}{:}}}}{1}
      {,}{{{\color{punct}{,}}}}{1}
      {\{}{{{\color{delim}{\{}}}}{1}
      {\}}{{{\color{delim}{\}}}}}{1}
      {[}{{{\color{delim}{[}}}}{1}
      {]}{{{\color{delim}{]}}}}{1},
}

% Define additional colors for JSON
\definecolor{numb}{rgb}{0.5,0,0}
\definecolor{punct}{rgb}{0,0,0}
\definecolor{delim}{rgb}{0.5,0.5,0.5}

% Header and footer
\pagestyle{fancy}
\fancyhf{}
\rhead{LangChain4j Integration Guide}
\lhead{Multi-Agent Insurance RAG System}
\rfoot{\thepage}

\title{
    \Huge\textbf{LangChain4j Integration Guide} \\
    \Large Multi-Agent Insurance Claim Processing System \\
    \large with Orchestrator-Workers Pattern
}
\author{
    \textbf{Technical Architecture Documentation} \\
    Spring Boot 3.3.5 + LangChain4j 0.34.0 + Spring AI 1.1.1
}
\date{\today}

\begin{document}

\maketitle
\tableofcontents
\newpage

\section{Executive Summary}

This document provides a comprehensive technical guide for the LangChain4j integration implemented within a Spring Boot application featuring a sophisticated multi-agent insurance claim processing system. The architecture leverages the \textbf{Orchestrator-Workers pattern} to coordinate specialized AI agents for automated claim classification, validation, and cost estimation.

\begin{tcolorbox}[colback=blue!5!white,colframe=blue!75!black,title=Key Achievements]
\begin{itemize}
    \item Successful dual-framework integration: LangChain4j + Spring AI
    \item Implementation of multi-agent insurance processing workflow
    \item RAG (Retrieval-Augmented Generation) for policy compliance validation
    \item Asynchronous orchestration with CompletableFuture
    \item Local LLM deployment with Ollama (llama3.2:latest)
\end{itemize}
\end{tcolorbox}

\section{System Architecture Overview}

\subsection{Technology Stack}

The integration combines multiple cutting-edge technologies:

\begin{center}
\begin{tabular}{|c|c|c|}
\hline
\textbf{Framework} & \textbf{Version} & \textbf{Purpose} \\
\hline
Spring Boot & 3.3.5 & Application Framework \\
LangChain4j & 0.34.0 & AI Orchestration \& RAG \\
Spring AI & 1.1.1 & Vector Store Integration \\
PostgreSQL & 16+ & Persistent Storage \\
PgVector & Latest & Vector Database Extension \\
Ollama & Latest & Local LLM Serving \\
\hline
\end{tabular}
\end{center}

\subsection{Multi-Agent Architecture}

The system implements the \textbf{Orchestrator-Workers pattern} with four specialized components:

\begin{enumerate}
    \item \textbf{AgentRouteur} - Classification agent for claim type identification
    \item \textbf{AgentValidateur} - RAG-powered compliance validation agent  
    \item \textbf{AgentEstimateur} - Multimodal cost estimation agent
    \item \textbf{OrchestratorMultiAgents} - Master coordinator for workflow management
\end{enumerate}

\section{LangChain4j Integration Implementation}

\subsection{Dependency Configuration}

The Maven configuration establishes LangChain4j integration alongside Spring AI:

\begin{lstlisting}[language=xml, caption=POM.xml LangChain4j Dependencies]
<properties>
    <langchain4j.version>0.34.0</langchain4j.version>
</properties>

<dependencies>
    <!-- LangChain4j Core -->
    <dependency>
        <groupId>dev.langchain4j</groupId>
        <artifactId>langchain4j</artifactId>
        <version>${langchain4j.version}</version>
    </dependency>
    
    <!-- Ollama Integration -->
    <dependency>
        <groupId>dev.langchain4j</groupId>
        <artifactId>langchain4j-ollama</artifactId>
        <version>${langchain4j.version}</version>
    </dependency>
    
    <!-- Local Embeddings -->
    <dependency>
        <groupId>dev.langchain4j</groupId>
        <artifactId>langchain4j-embeddings-all-minilm-l6-v2</artifactId>
        <version>${langchain4j.version}</version>
    </dependency>
    
    <!-- PDF Document Processing -->
    <dependency>
        <groupId>dev.langchain4j</groupId>
        <artifactId>langchain4j-document-parser-apache-tika</artifactId>
        <version>${langchain4j.version}</version>
    </dependency>
</dependencies>
\end{lstlisting}

\subsection{Configuration Class Architecture}

The \texttt{OllamaConfig} class manages bean instantiation and framework coexistence:

\begin{lstlisting}[language=java, caption=OllamaConfig.java - Dual Framework Configuration]
@Configuration
public class OllamaConfig {

    @Bean
    @Qualifier("langchain4j")
    public ChatLanguageModel chatLanguageModel() {
        return OllamaChatModel.builder()
                .baseUrl("http://localhost:11434")
                .modelName("llama3.2:latest")
                .temperature(0.3)
                .build();
    }

    @Bean  
    @Qualifier("langchain4j")
    public EmbeddingModel embeddingModel() {
        return OllamaEmbeddingModel.builder()
                .baseUrl("http://localhost:11434")
                .modelName("nomic-embed-text:latest")
                .build();
    }

    @Bean("langchain4jChatMemory")
    public ChatMemory langchain4jChatMemory() {
        return MessageWindowChatMemory.withMaxMessages(20);
    }

    // RAG Components
    @Bean
    public InMemoryEmbeddingStore<TextSegment> embeddingStore() {
        return new InMemoryEmbeddingStore<>();
    }

    @Bean
    public ContentRetriever contentRetriever(
            EmbeddingModel embeddingModel,
            EmbeddingStore<TextSegment> embeddingStore) {
        return EmbeddingStoreContentRetriever.builder()
                .embeddingStore(embeddingStore)
                .embeddingModel(embeddingModel)
                .maxResults(3)
                .minScore(0.7)
                .build();
    }
}
\end{lstlisting}

\section{Multi-Agent System Implementation}

\subsection{Agent Routeur - Classification Engine}

The classification agent uses natural language processing to categorize insurance claims:

\begin{lstlisting}[language=java, caption=AgentRouteur.java - Intelligent Classification]
@Service
public class AgentRouteur {
    
    private final ChatLanguageModel chatModel;
    
    public SinistreType classifierTypeSinistre(DemandeTraitement demande) {
        String prompt = """
            Analysez cette description de sinistre et classifiez-la:
            "%s"
            
            Types disponibles:
            - ACCIDENT_AUTO: accidents de v\'{e}hicules
            - DEGAT_EAU: fuites, inondations, d\'{e}g\^{a}ts des eaux
            - INCENDIE: feux, explosions, fum\'{e}e
            - VOL: cambriolage, vol \`{a} l'arrach\'{e}
            - VANDALISME: d\'{e}gradations volontaires
            - BRIS_GLACE: pare-brise, vitres, miroirs
            - CATASTROPHE_NATURELLE: temp\^{e}te, gr\^{e}le, s\'{e}isme
            - AUTRES: tout autre type de sinistre
            
            R\'{e}pondez uniquement par le nom exact du type.
            """.formatted(demande.getContenuDemande());

        String reponse = chatModel.generate(prompt);
        
        // Parse and validate response with fallback logic
        return parseTypeSinistre(reponse, demande.getContenuDemande());
    }
}
\end{lstlisting}

\subsection{Agent Validateur - RAG-Powered Compliance}

The validation agent integrates RAG capabilities for policy compliance checking:

\begin{lstlisting}[language=java, caption=AgentValidateur.java - RAG Integration]
@Service
public class AgentValidateur {
    
    private final ChatLanguageModel chatModel;
    private final LangChain4jRagService ragService;
    
    public boolean validerConformite(DemandeTraitement demande) {
        // Query policy database using RAG
        String questionPolitique = String.format(
            "Quelles sont les conditions de couverture pour %s? " +
            "Y a-t-il des exclusions pour: %s",
            demande.getTypeSinistre().getDescription(),
            demande.getContenuDemande()
        );
        
        String informationsPolitique = ragService.obtenirInformations(questionPolitique);
        
        String prompt = """
            Validez la conformit\'{e} de ce sinistre avec la police d'assurance:
            
            Type: %s
            Description: "%s"
            Informations de la police: %s
            
            V\'{e}rifiez:
            1. Le type de sinistre est-il couvert?
            2. Les circonstances respectent-elles les conditions?
            3. Aucune exclusion ne s'applique-t-elle?
            4. Les d\'{e}lais de d\'{e}claration sont-ils respect\'{e}s?
            5. Les informations fournies sont-elles suffisantes?
            
            R\'{e}pondez par CONFORME ou NON\_CONFORME avec justification d\'{e}taill\'{e}e.
            """.formatted(
                demande.getTypeSinistre().getDescription(),
                demande.getContenuDemande(),
                informationsPolitique
            );

        String reponse = chatModel.generate(prompt);
        return !reponse.toUpperCase().contains("NON_CONFORME");
    }
}
\end{lstlisting}

\subsection{Agent Estimateur - Multimodal Analysis}

The estimation agent performs sophisticated damage assessment with simulated multimodal capabilities:

\begin{lstlisting}[language=java, caption=AgentEstimateur.java - Cost Analysis Engine]
@Service 
public class AgentEstimateur {
    
    public double estimerCout(DemandeTraitement demande) {
        String analysePhotos = analyserPhotos(demande.getPhotosUrls());
        
        String prompt = """
            En tant qu'expert estimateur d'assurance, calculez une estimation:
            
            Type: %s
            Description: "%s" 
            Analyse photos: %s
            
            Basez votre estimation sur:
            1. Type et \'{e}tendue des dommages
            2. Co\^{u}ts moyens de r\'{e}paration/remplacement
            3. Main d'\oe{}uvre n\'{e}cessaire
            4. Pi\`{e}ces et mat\'{e}riaux requis
            5. Co\^{u}ts indirects (expertise, location, etc.)
            
            Format: [MONTANT] | [JUSTIFICATION] | [CONFIANCE]
            """.formatted(
                demande.getTypeSinistre().getDescription(),
                demande.getContenuDemande(),
                analysePhotos
            );

        String reponse = chatModel.generate(prompt);
        return extraireMontant(reponse);
    }
    
    private String analyserPhotos(List<String> photosUrls) {
        // Simulated multimodal analysis
        // In production: integrate GPT-4V or Claude 3 for real image analysis
        StringBuilder analyse = new StringBuilder();
        // ... sophisticated photo analysis logic
        return analyse.toString();
    }
}
\end{lstlisting}

\subsection{Orchestrator - Asynchronous Workflow Coordination}

The orchestrator manages the complete claim processing workflow using asynchronous programming:

\begin{lstlisting}[language=java, caption=OrchestratorMultiAgents.java - Workflow Management]
@Service
public class OrchestratorMultiAgents {
    
    private final ExecutorService executorService = 
        Executors.newFixedThreadPool(10);
    
    public CompletableFuture<DemandeTraitement> traiterDemandeComplete(
            DemandeTraitement demande) {
        
        return CompletableFuture
            .supplyAsync(() -> {
                // Phase 1: Classification
                SinistreType type = agentRouteur.classifierTypeSinistre(demande);
                demande.setTypeSinistre(type);
                return demande;
            }, executorService)
            .thenCompose(demandeAvecType -> 
                CompletableFuture.supplyAsync(() -> {
                    // Phase 2: Validation
                    boolean conforme = agentValidateur.validerConformite(demandeAvecType);
                    demandeAvecType.setConformite(conforme);
                    
                    if (!conforme) {
                        demandeAvecType.setStatut(StatutTraitement.REJETE);
                        return demandeAvecType;
                    }
                    
                    return demandeAvecType;
                }, executorService)
            )
            .thenCompose(demandeValidee -> {
                if (demandeValidee.getStatut() == StatutTraitement.REJETE) {
                    return CompletableFuture.completedFuture(demandeValidee);
                }
                
                // Phase 3: Cost Estimation
                return CompletableFuture.supplyAsync(() -> {
                    double cout = agentEstimateur.estimerCout(demandeValidee);
                    demandeValidee.setEstimationCout(cout);
                    demandeValidee.setStatut(StatutTraitement.ESTIME);
                    return demandeValidee;
                }, executorService);
            });
    }
}
\end{lstlisting}

\section{RAG (Retrieval-Augmented Generation) Implementation}

\subsection{LangChain4j RAG Service}

The RAG service provides sophisticated document retrieval and generation capabilities:

\begin{lstlisting}[language=java, caption=LangChain4jRagService.java - RAG Pipeline]
@Service
public class LangChain4jRagService {
    
    private final ContentRetriever contentRetriever;
    private final ChatLanguageModel chatModel;
    
    public String obtenirInformations(String question) {
        // Retrieve relevant documents
        List<Content> contenusPertinents = contentRetriever.retrieve(
            Query.from(question)
        );
        
        if (contenusPertinents.isEmpty()) {
            return "Aucune information spécifique trouvée dans la base documentaire.";
        }
        
        // Build context from retrieved documents
        String contexte = contenusPertinents.stream()
            .map(Content::textSegment)
            .map(TextSegment::text)
            .collect(Collectors.joining("\\n\\n"));
            
        String prompt = """
            Bas\'{e} uniquement sur ces informations documentaires:
            %s
            
            Question: %s
            
            R\'{e}pondez en utilisant exclusivement les informations fournies.
            Si l'information n'est pas disponible, indiquez-le clairement.
            """.formatted(contexte, question);
            
        return chatModel.generate(prompt);
    }
}
\end{lstlisting}

\subsection{Document Ingestion Pipeline}

The document processing pipeline transforms PDFs into searchable vector embeddings:

\begin{lstlisting}[language=java, caption=VectorDocumentIngestionService.java - Document Processing]
@Service  
public class VectorDocumentIngestionService {
    
    public void ingestDocumentToLangChain4j(String filePath) {
        try {
            // Load and parse document
            Document document = loadDocument(toPath(filePath), 
                new ApacheTikaDocumentParser());
                
            // Split into chunks with overlap
            DocumentSplitter splitter = DocumentSplitters
                .recursive(chunkSize, chunkOverlap);
            List<TextSegment> segments = splitter.split(document);
            
            // Generate embeddings and store
            List<Embedding> embeddings = embeddingModel.embedAll(segments);
            embeddingStore.addAll(embeddings, segments);
            
            logger.info("Successfully ingested {} segments into LangChain4j vector store", 
                segments.size());
                
        } catch (Exception e) {
            logger.error("Error during LangChain4j document ingestion", e);
        }
    }
}
\end{lstlisting}

\section{API Endpoints and Usage}

\subsection{Multi-Agent Processing Endpoints}

The system exposes RESTful endpoints for comprehensive claim processing:

\begin{lstlisting}[language=java, caption=MultiAgentsController.java - REST API]
@RestController
@RequestMapping("/api/agents")
public class MultiAgentsController {
    
    @GetMapping("/status")
    public ResponseEntity<Map<String, Object>> getSystemStatus() {
        return ResponseEntity.ok(Map.of(
            "service", "Multi-Agents Orchestrator",
            "status", "Op\\'erationnel", 
            "agents", List.of("Routeur", "Validateur", "Estimateur"),
            "patterns", "Orchestrator-Workers",
            "capabilities", List.of(
                "Classification automatique des sinistres",
                "Validation de conformit\\'e via RAG", 
                "Estimation multimodale des co\\'uts",
                "Rapports d\\'etaill\\'es",
                "Traitement asynchrone"
            )
        ));
    }
    
    @PostMapping("/traiter-sinistre")
    public CompletableFuture<ResponseEntity<Map<String, Object>>> 
            traiterSinistre(@RequestBody Map<String, Object> requestData) {
        
        String email = (String) requestData.get("email");
        String description = (String) requestData.get("description"); 
        List<String> photos = (List<String>) requestData.getOrDefault("photos", List.of());
        
        DemandeTraitement demande = new DemandeTraitement(email, description, photos);
        
        return orchestrateur.traiterDemandeComplete(demande)
            .thenApply(demandeTraitee -> ResponseEntity.ok(Map.of(
                "success", true,
                "demandeId", demandeTraitee.getId(),
                "statut", demandeTraitee.getStatut().name(),
                "typeSinistre", demandeTraitee.getTypeSinistre().getDescription(),
                "conforme", demandeTraitee.isConformite(),
                "estimationCout", demandeTraitee.getEstimationCout()
            )));
    }
}
\end{lstlisting}

\subsection{Usage Examples}

\subsubsection{System Status Check}

\begin{lstlisting}[language=bash, caption=Check System Status]
curl -X GET "http://localhost:8080/api/agents/status"
\end{lstlisting}

Expected Response:
\begin{lstlisting}[language=json]
{
  "service": "Multi-Agents Orchestrator",
  "status": "Op\\'erationnel",
  "agents": ["Routeur", "Validateur", "Estimateur"],
  "patterns": "Orchestrator-Workers",
  "capabilities": [
    "Classification automatique des sinistres",
    "Validation de conformit\\'e via RAG",
    "Estimation multimodale des co\\'uts"
  ]
}
\end{lstlisting}

\subsubsection{Insurance Claim Processing}

\begin{lstlisting}[language=bash, caption=Process Insurance Claim]
curl -X POST "http://localhost:8080/api/agents/traiter-sinistre" \
  -H "Content-Type: application/json" \
  -d '{
    "email": "client@example.com",
    "description": "Accident de voiture sur autoroute A7, collision frontale",
    "photos": ["collision_avant.jpg", "degats_droite.jpg"]
  }'
\end{lstlisting}

Expected Response:
\begin{lstlisting}[language=json]
{
  "success": true,
  "demandeId": "eff094e0-0569-4948-bc04-2b769cd9a52f", 
  "statut": "ESTIME",
  "typeSinistre": "Accident automobile",
  "conforme": true,
  "estimationCout": 4250.00
}
\end{lstlisting}

\section{Configuration and Deployment}

\subsection{Application Properties}

The application requires specific configuration for dual-framework operation:

\begin{lstlisting}[caption=application.properties - Complete Configuration]
# Application
spring.application.name=demo

# Database Configuration  
spring.datasource.url=jdbc:postgresql://localhost:5432/testrag
spring.datasource.username=postgres
spring.datasource.password=postgres
spring.jpa.hibernate.ddl-auto=update

# Spring AI Configuration
spring.ai.ollama.base-url=http://localhost:11434
spring.ai.ollama.chat.options.model=llama3.2:latest
spring.ai.ollama.embedding.options.model=nomic-embed-text:latest
spring.ai.vectorstore.pgvector.initialize-schema=true

# LangChain4j Configuration
langchain4j.ollama.chat-model.base-url=http://localhost:11434
langchain4j.ollama.chat-model.model-name=llama3.2:latest
langchain4j.ollama.embedding-model.base-url=http://localhost:11434
langchain4j.ollama.embedding-model.model-name=nomic-embed-text:latest

# Document Processing
app.pdf.chunk-size=500
app.pdf.chunk-overlap=100
\end{lstlisting}

\subsection{Docker Deployment}

For containerized deployment, the following Docker Compose configuration ensures all services are properly coordinated:

\begin{lstlisting}[language=yaml, caption=docker-compose.yml]
version: '3.8'
services:
  postgres:
    image: pgvector/pgvector:pg16
    environment:
      POSTGRES_DB: testrag
      POSTGRES_USER: postgres  
      POSTGRES_PASSWORD: postgres
    ports:
      - "5432:5432"
    volumes:
      - postgres_data:/var/lib/postgresql/data
      
  ollama:
    image: ollama/ollama:latest
    ports:
      - "11434:11434"
    volumes:
      - ollama_data:/root/.ollama
    command: >
      sh -c "ollama serve & 
             sleep 10 && 
             ollama pull llama3.2:latest &&
             ollama pull nomic-embed-text:latest &&
             wait"

volumes:
  postgres_data:
  ollama_data:
\end{lstlisting}

\section{Performance Optimization and Best Practices}

\subsection{Asynchronous Processing Strategies}

The multi-agent system implements several performance optimization techniques:

\begin{itemize}
    \item \textbf{CompletableFuture Chains}: Non-blocking sequential agent execution
    \item \textbf{Thread Pool Management}: Dedicated executor service for agent tasks  
    \item \textbf{Caching Mechanisms}: RAG result caching for repeated policy queries
    \item \textbf{Connection Pooling}: Optimized database and LLM API connections
\end{itemize}

\subsection{Error Handling and Resilience}

\begin{tcolorbox}[colback=red!5!white,colframe=red!75!black,title=Error Handling Strategy]
\begin{enumerate}
    \item \textbf{Circuit Breaker Pattern}: Prevents cascading failures in agent communication
    \item \textbf{Retry Logic}: Automatic retry with exponential backoff for transient failures
    \item \textbf{Fallback Mechanisms}: Default behaviors when AI services are unavailable
    \item \textbf{Graceful Degradation}: Partial functionality maintenance during service outages
\end{enumerate}
\end{tcolorbox}

\subsection{Monitoring and Observability}

The system includes comprehensive monitoring capabilities:

\begin{lstlisting}[language=java, caption=Monitoring Integration]
@Component
public class AgentMetrics {
    
    private final MeterRegistry meterRegistry;
    private final Counter claimsProcessed;
    private final Timer processingTime;
    
    public AgentMetrics(MeterRegistry meterRegistry) {
        this.meterRegistry = meterRegistry;
        this.claimsProcessed = Counter.builder("claims.processed")
            .description("Total claims processed")
            .register(meterRegistry);
        this.processingTime = Timer.builder("claims.processing.time")
            .description("Claim processing duration")
            .register(meterRegistry);
    }
}
\end{lstlisting}

\section{Future Enhancements and Roadmap}

\subsection{Planned Improvements}

\begin{enumerate}
    \item \textbf{Real Multimodal Integration}: Replace photo simulation with actual GPT-4V or Claude 3 Vision
    \item \textbf{Enhanced RAG Pipeline}: Implement hybrid search (dense + sparse retrieval)
    \item \textbf{Agent Learning}: Add feedback loops for continuous agent improvement
    \item \textbf{Blockchain Integration}: Immutable claim processing audit trails
    \item \textbf{Real-time Streaming}: WebSocket-based real-time claim status updates
\end{enumerate}

\subsection{Scalability Considerations}

For enterprise deployment, consider the following architectural enhancements:

\begin{itemize}
    \item \textbf{Microservices Decomposition}: Split agents into independent services
    \item \textbf{Event-Driven Architecture}: Replace synchronous calls with event streaming
    \item \textbf{Distributed Caching}: Redis cluster for cross-instance data sharing
    \item \textbf{Load Balancing}: Multiple Ollama instances behind load balancer
    \item \textbf{Database Sharding}: Partition vector data across multiple PostgreSQL instances
\end{itemize}

\section{Troubleshooting Guide}

\subsection{Common Issues and Solutions}

\begin{tcolorbox}[colback=yellow!5!white,colframe=yellow!75!black,title=Troubleshooting Checklist]
\textbf{Issue}: Bean naming conflicts between Spring AI and LangChain4j
\textbf{Solution}: Use \texttt{@Qualifier} annotations and unique bean names

\textbf{Issue}: Ollama connection timeout
\textbf{Solution}: Verify Ollama service status with \texttt{curl http://localhost:11434/api/tags}

\textbf{Issue}: Vector store initialization errors  
\textbf{Solution}: Ensure PostgreSQL has pgvector extension: \texttt{CREATE EXTENSION vector;}

\textbf{Issue}: Out of memory during document ingestion
\textbf{Solution}: Increase JVM heap size: \texttt{-Xmx4g -Xms2g}
\end{tcolorbox}

\subsection{Performance Tuning}

For optimal system performance:

\begin{lstlisting}[caption=JVM Optimization Parameters]
# Production JVM Settings
-Xmx8g
-Xms4g  
-XX:+UseG1GC
-XX:MaxGCPauseMillis=200
-XX:+UnlockExperimentalVMOptions
-XX:+UseStringDeduplication
\end{lstlisting}

\section{Conclusion}

The LangChain4j integration successfully demonstrates advanced AI orchestration capabilities within a Spring Boot ecosystem. The multi-agent insurance claim processing system showcases:

\begin{itemize}
    \item \textbf{Framework Coexistence}: Seamless operation of LangChain4j alongside Spring AI
    \item \textbf{Sophisticated AI Workflows}: Complex multi-step claim processing with specialized agents
    \item \textbf{Production-Ready Architecture}: Asynchronous processing, error handling, and monitoring
    \item \textbf{Scalable Design}: Ready for enterprise deployment with microservices evolution
\end{itemize}

This implementation serves as a blueprint for building sophisticated AI-powered business applications that combine the strengths of multiple AI frameworks while maintaining clean architecture principles and operational excellence.

\vfill
\textit{This documentation represents a comprehensive technical guide for the LangChain4j integration implementation. For updates and additional resources, refer to the project repository and official LangChain4j documentation.}

\end{document}
